\section{Design}
\label{sec:design}

\subsection{Overview}

\Cref{fig:architecture} shows the architecture.  Each managed cgroup owns a
\emph{phase sketch}, a fixed-size table relating execution signatures to page
regions.  Hardware samples flow through the kernel perf ring buffer; a worker
coalesces them into an epoch signature.  On a confident signature match, the
predictor proposes regions to the promotion engine.  A global credit controller
limits copy traffic, and per-cgroup leases bound occupancy.  Conventional
access-based promotion remains active for unrecognized phases.

\begin{figure*}[t]
  \centering
  \begin{tikzpicture}[
    node distance=0.72cm and 0.72cm,
    block/.style={draw=fluxblue,thick,rounded corners=2pt,fill=fluxblue!7,
                  minimum width=2.55cm,minimum height=0.88cm,align=center,
                  font=\small},
    control/.style={draw=fluxorange,thick,rounded corners=2pt,
                    fill=fluxorange!10,minimum width=2.55cm,
                    minimum height=0.88cm,align=center,font=\small},
    mem/.style={draw=fluxteal,thick,rounded corners=2pt,fill=fluxteal!10,
                minimum width=2.30cm,minimum height=0.95cm,align=center,
                font=\small},
    arrow/.style={-Latex,thick,draw=fluxgray},
    edgelabel/.style={font=\scriptsize,text=fluxgray,fill=white,inner sep=1pt}
  ]
    \node[block] (perf) {IP + address\\samples};
    \node[block,right=of perf] (sig) {epoch signature\\and phase sketch};
    \node[control,right=of sig] (pred) {confidence and\\benefit estimator};
    \node[control,right=of pred] (credit) {migration-credit\\controller};
    \node[block,below=1.05cm of sig] (fallback) {reactive hot-page\\scanner};
    \node[mem,right=0.60cm of fallback] (dram) {local DRAM\\\fast};
    \node[mem,right=1.20cm of dram] (cxl) {CXL memory\\\slow};
    \node[control,right=0.60cm of cxl] (lease) {lease expiry and\\deadline-aware reclaim};

    \draw[arrow] (perf) -- (sig);
    \draw[arrow] (sig) -- (pred);
    \draw[arrow] (pred) -- (credit);
    \draw[arrow] (credit) |- (cxl) node[edgelabel,pos=0.28] {admit};
    \draw[arrow] (cxl) -- (dram) node[edgelabel,midway,below] {promote};
    \draw[arrow] (dram) to[bend left=32]
      node[edgelabel,midway,above] {demote} (cxl);
    \draw[arrow] (perf) |- (fallback);
    \draw[arrow] (fallback) -- (dram);
    \draw[arrow] (lease) -- (cxl);
    \draw[arrow,dashed] (pred.south) -- ++(0,-0.35cm) -| (lease.north);

    \node[draw=black!25,dashed,rounded corners=3pt,
          fit=(sig)(pred)(credit)(fallback)(dram)(cxl)(lease),
          inner sep=8pt,label={[font=\small\bfseries,text=fluxblue]above:\system}]
          {};
  \end{tikzpicture}
  \caption{\system data and control paths.  Prediction augments rather than
    replaces reactive hot-page discovery.  Credits bound migration bandwidth;
    leases bound speculative occupancy.}
  \Description{A block diagram shows instruction-pointer and address samples
    entering an epoch signature and phase sketch, then a confidence estimator
    and migration-credit controller. Predicted and reactive paths move pages
    between local DRAM and CXL memory; leases govern reclaim.}
  \label{fig:architecture}
\end{figure*}

\subsection{Phase signatures}

An epoch contains $n$ sampled tuples $(i,a,r)$: instruction pointer $i$,
virtual address $a$, and read/write bit $r$.  We discard kernel addresses and
hash the remaining instruction pointers into a 128-bit Bloom-style signature
$S$.  We deliberately omit addresses from $S$ so that address-space layout
randomization and allocator variation do not prevent a match.  The signature
distance is normalized Hamming distance,

\begin{equation}
  d(S_x,S_y)=\frac{\operatorname{popcount}(S_x\oplus S_y)}
                   {\operatorname{popcount}(S_x\lor S_y)+1}.
  \label{eq:distance}
\end{equation}

The phase table is a 64-entry, four-way set-associative structure per cgroup.
Each entry stores a signature, an exponentially weighted transition count,
the eight most frequently touched 2\,MiB regions in the next two epochs, and a
4-bit confidence counter.  Regions identify sets of base pages; \system never
promotes a huge page unless all constituent base pages are resident and the
copy engine can migrate it without blocking.

On each epoch boundary, the nearest entry with $d<0.18$ receives an update.  A
miss installs a new entry only after the signature appears twice, filtering
one-off control paths.  Confidence increments when at least half of the
predicted regions enter the next epoch's top set and decrements otherwise.
The fixed-size table is essential: at 416 bytes per entry, predictor metadata
is 26\,KiB per actively managed cgroup, independent of its address-space size.

\subsection{Benefit and confidence}

A signature match does not automatically trigger migration.  For candidate
region $q$, the predictor estimates benefit

\begin{equation}
  B(q)=P(q\mid S)\,N_q\,(L_{\slow}-L_{\fast})-C_{\mathrm{copy}}(q),
  \label{eq:benefit}
\end{equation}

where $P(q\mid S)$ is the entry's hit ratio, $N_q$ is its decayed access count,
and $L_{\slow}-L_{\fast}$ is the measured latency gap.  The copy cost combines
expected bandwidth occupancy and the cost of displacing the coldest local
pages.  A candidate is eligible only if $B(q)>0$, confidence is at least 9 of
15, and the destination page is not already resident in \fast.

The estimates need not be cycle-accurate.  Their purpose is to rank candidates
and reject clearly harmful moves.  The credit controller supplies the hard
safety bound described next.

\subsection{Credit-bounded promotion}

Copy engines share memory channels with foreground loads.  \system models each
4\,KiB migration as one credit.  Every 100\,\us, the system replenishes at rate

\begin{equation}
  R = \frac{\beta\,BW_{\mathrm{idle}}\,\Delta t}{4096},
  \label{eq:credits}
\end{equation}

where $BW_{\mathrm{idle}}$ is currently unused channel bandwidth and
$\beta=0.35$ caps the fraction made available to prediction.  Credits can
accumulate for at most 2\,ms, preventing an idle period from authorizing a later
burst.  Each cgroup receives a weighted share, but unused shares are
work-conservingly borrowed.  Reactive promotions use a distinct reserve equal
to 10\% of copy capacity, ensuring that bad prediction cannot disable the
fallback path.

\begin{lstlisting}[float=t,caption={Promotion at an epoch boundary.},
                   label={lst:promotion}]
procedure ON_EPOCH(cgroup g, signature S):
  e <- nearest(g.phase_table, S)
  if e = none or e.confidence < 9 then return
  for each q in rank_by_benefit(e.regions) do
    if benefit(q) <= 0 then continue
    pages <- cold_pages(q, g.fast_residency)
    cost <- copy_credits(pages)
    if cost <= g.credits and lease_slot(g) then
      migrate_async(pages, M_fast)
      g.credits <- g.credits - cost
      attach_lease(pages, now + lease_length(e))
\end{lstlisting}

\subsection{Leases and deadline-aware reclaim}

Promoted pages receive a lease of two predicted phase durations, clamped to
1--20\,ms.  During the lease, ordinary background reclaim may select the page
only if \fast falls below its 3\% emergency reserve.  A sampled access renews
the lease; expiry merely removes protection and does not force a migration.
Pages then compete under the regular multigenerational LRU policy.

Each cgroup additionally has a local-residency ceiling.  A prediction that
would exceed the ceiling must nominate expired pages from the same cgroup
before borrowing globally.  This rule prevents a noisy tenant from converting
shared \fast into a private cache.  Under acute pressure, the emergency path
ignores every lease and reclaims in ordinary LRU order, so \system cannot cause
an out-of-memory failure by pinning data.

\subsection{Failure handling}

Prediction is an optimization, not a correctness requirement.  Process exit,
unmap, and cgroup deletion cancel queued copies through existing migration
reference counts.  A copy racing with a write uses Linux's normal page
migration protocol and page-table locks.  If the destination fills, migration
returns \texttt{-ENOMEM}; the page remains mapped in \slow.  Sampling loss,
perf-ring overflow, or confidence decay disables prediction for that cgroup
until two stable signatures recur.  Thus every failure returns to the baseline
reactive policy.
